% ===================================================================== Chapter 9
\chapter{CFD Mesh Independence}\label{ch:mi}

\section{Purpose}
The mesh study estimates how much of each baseline result is discretisation error, and whether the medium mesh is adequate as the thermal
input to the structural analysis. The coarse and fine meshes were solved with physics, boundary conditions, solution staging and convergence
criteria identical to the medium baseline. Both converged at iteration 600 and held at 700; the fine solution conserves mass and energy to
$3.9\times10^{-12}$\,\% and $6.2\times10^{-11}$\,\%. Table~\ref{tab:mi} lists the results and Figure~\ref{fig:mi} plots four of them.

%% generated by gen_tables.py - RE-ANALYSIS 2026
\begin{table}[htbp]
\centering
\footnotesize
\caption{CFD mesh-independence study: three meshes and Richardson extrapolation}
\label{tab:mi}
\begin{tabular}{lrrrrl}
\toprule
Quantity & Coarse & Medium & Fine & Extrapolated & Status \\
\midrule
Cells & 51,840 & 159,840 & 500,580 & -- & -- \\
$\Delta p$ [Pa] & 437.09 & 438.13 & 439.12 & 441.81$^a$ & B \\
$T_{out}$ [K] & 369.12 & 368.93 & 368.85 & -- & D (geometry) \\
$Q$ [W] & 602.217 & 602.755 & 602.994 & 603.186$^b$ & D (geometry) \\
$T_{max}$ solid [K] & 565.40 & 562.58 & 560.83 & 558.13 & B \\
$\Delta T_{wall}$ mid-span [K] & 7.554 & 7.581 & 7.597 & 7.620 & B \\
$Nu_{fd}$ & 53.38 & 54.17 & 54.66 & 55.43 & B \\
$f_{fd}$ & 0.02134 & 0.02163 & 0.02181 & 0.02212 & B \\
$y^+$ maximum & 0.522 & 0.585 & 0.653 & -- & B / E \\
\bottomrule
\end{tabular}
\par\smallskip\parbox{\linewidth}{\footnotesize\raggedright Source: \texttt{MASTER\_PROJECT\_DATA.csv} rows M033--M070. Status: B approximately convergent; D affected by the polygon geometry; E inconclusive (sampling position). $^a$ geometry-adjusted estimate. $^b$ exact circle value recovered by extrapolating the raw values (check of the procedure).}
\end{table}


\section{Coarse Mesh}
The coarse mesh (51,840~cells) gives a maximum solid temperature of 565.40~K, 7.3~K above the extrapolated value; from coarse to medium the Nusselt number
changes by +1.48\,\%. It was used only to complete the three-level family and was rejected for downstream use.

\section{Medium Mesh}
The medium mesh (159,840~cells) is the official baseline. Its maximum solid temperature is 4.4~K above the extrapolated value and its
volume-mean solid temperature 3.1~K above; the mid-span through-wall temperature difference is 0.04~K low. For the structural analysis these
errors are small and act in the conservative direction (slightly hotter wall, slightly larger restrained force).

\section{Fine Mesh}
The fine mesh (500,580~cells) is the verification reference. From medium to fine the pressure drop changes by +0.23\,\%, the maximum solid
temperature by $-$1.75~K ($-$0.66\,\% of the temperature rise), the Nusselt number by +0.91\,\%, the friction factor by +0.85\,\% and the
through-wall temperature difference by +0.21\,\%. Its maximum solid temperature, 560.83~K, is quoted wherever the fine result is referred to.

\begin{figure}[htbp]
\centering
\begin{subfigure}[t]{0.49\linewidth}\centering\includegraphics[width=\linewidth,height=52mm,keepaspectratio]{MI01_cells_vs_pressure_drop.png}\caption{Pressure drop [Pa]}\end{subfigure}\hfill
\begin{subfigure}[t]{0.49\linewidth}\centering\includegraphics[width=\linewidth,height=52mm,keepaspectratio]{MI04_cells_vs_max_solid_temperature.png}\caption{Maximum solid temperature [K]}\end{subfigure}\\[4pt]
\begin{subfigure}[t]{0.49\linewidth}\centering\includegraphics[width=\linewidth,height=52mm,keepaspectratio]{MI06_cells_vs_nusselt.png}\caption{Fully developed Nusselt number}\end{subfigure}\hfill
\begin{subfigure}[t]{0.49\linewidth}\centering\includegraphics[width=\linewidth,height=52mm,keepaspectratio]{MI07_cells_vs_friction_factor.png}\caption{Fully developed Darcy friction factor}\end{subfigure}
\caption[Mesh-independence study]{Mesh-independence study: raw CFD results on the three meshes, geometry-adjusted values, Richardson extrapolation and GCI$_{fine}$ error bars.}
\label{fig:mi}
\src{\provcfd{} In (b) the extrapolated line belongs to the bulk-corrected series ($p$ = 1.26, 558.06~K); the raw series gives $p$ = 1.31 and 558.13~K (Table~\ref{tab:gci}).}
\end{figure}

\section{Geometry/Faceting Effect}
The three meshes approximate the circular sections by inscribed polygons with 32, 48 and 72 sides. The flow area, heated area and therefore
the mass flow, heat input and outlet temperature change with the number of sides by exactly known ratios. These changes are geometry, not
discretisation error, and were separated exactly before any convergence assessment (decision D-035). Corrected to the true circle, the mass
flow is 8.68694~g/s and the heat input 603.186~W on all three meshes; Richardson extrapolation of the raw values recovers the exact circle
values to within 0.004\,\%, which checks the procedure. Of the raw changes, nearly all of those in mass flow, heat rate and outlet temperature are geometric,
only 5--7\,\% of the change in maximum solid temperature, and a negligible part of the changes in Nusselt number and friction factor.


\section{Pressure-Drop Convergence}
The raw pressure drop increases monotonically (437.09, 438.13, 439.12~Pa), but the geometric part of the change opposes the resolution trend,
so the raw values give no credible order. The geometry-adjusted estimate has an apparent order of 1.05 and extrapolates to 441.8~Pa with a
fine-mesh GCI of 0.89\,\% (3.9~Pa). The medium value is 0.8\,\% below that estimate.

\section{Thermal Convergence}
The maximum solid temperature decreases monotonically (565.40, 562.58, 560.83~K) with an apparent order of 1.31, extrapolating to 558.13~K
with GCI$_{fine}$ = 3.4~K, 1.3\,\% of the temperature rise (Table~\ref{tab:gci}). Refinement raises the heat-transfer coefficient slightly
and therefore lowers the wall temperature, which moves the CFD further from the analytical 581.71~K. The mid-span through-wall temperature
difference converges to 7.620~K with GCI$_{fine}$ = 0.029~K.

%% generated by gen_tables.py - RE-ANALYSIS 2026
\begin{table}[htbp]
\centering
\footnotesize
\caption{Convergence indicators (Celik et al.\ procedure, $r_{21}$ = 1.4631, $r_{32}$ = 1.4555)}
\label{tab:gci}
\begin{tabular}{lrrrr}
\toprule
Quantity & Ratio $R$ & Apparent order $p$ & Extrapolated & GCI$_{fine}$ \\
\midrule
$\Delta p$, geometry-adjusted & 0.68 & 1.05 & 441.8 Pa & 0.89\,\% (3.9 Pa) \\
$T_{max}$ solid & 0.62 & 1.31 & 558.1 K & 3.4 K (1.3\,\% of rise) \\
$\Delta T_{wall}$ mid-span & 0.60 & 1.37 & 7.620 K & 0.029 K \\
$Nu_{fd}$ & 0.62 & 1.30 & 55.43 & 1.8\,\% \\
$f_{fd}$ & 0.64 & 1.23 & 0.02212 & 1.8\,\% \\
\bottomrule
\end{tabular}
\par\smallskip\parbox{\linewidth}{\footnotesize\raggedright Source: \texttt{PROJECT\_STATE.md} \S12.4 and \texttt{09\_Mesh\_Independence/Mesh\_Independence\_Report.md} \cite{celik2008}. Convergence is monotonic but the apparent order is below the formal value of 2.}
\end{table}


\section{Nusselt Number Convergence}
The fully developed Nusselt number increases monotonically (53.38, 54.17, 54.66) with an apparent order of 1.30 and extrapolates to 55.43
(GCI 1.8\,\%). Every value lies inside the expected band of 44--67. Convergence is monotonic for all quantities, but the apparent orders of
1.2--1.6 are below the formal second order, because the first layer is identical on all meshes and the directional refinement ratios are not
equal. The status of every flow and thermal quantity is therefore ``approximately convergent''; strict mesh independence is not claimed.

\section{Friction-Factor Behavior}
The friction factor rises with refinement (0.02134, 0.02163, 0.02181) towards an extrapolated 0.02212, still 8.4\,\% below the Petukhov
value. An isothermal diagnostic run on the medium mesh (energy equation off) gives $f$ = 0.02304, only 2.6\,\% below Petukhov
(Figure~\ref{fig:mifric}). The medium-mesh gap of $-$10.4\,\% decomposes exactly into a Reynolds-number basis factor of 1.011, an isothermal SST
factor of 0.974 and a heating factor of 0.910. The difference is therefore mainly a gas-heating (variable-property) effect, not a mesh defect;
about five percentage points remain unexplained by the standard $(T_w/T_b)^{-0.1}$ heating correction (finding F-034).

\begin{figure}[htbp]
\centering
\includegraphics[width=0.68\linewidth]{MI11_F029_friction_evidence_nt.png}
\caption{Local Darcy friction factor on the three heated meshes and in the isothermal diagnostic, compared with Petukhov with and without the heating correction.}
\label{fig:mifric}
\src{\provcfd}
\end{figure}

\section{Final Mesh Selection}
The medium mesh was retained for all downstream work, with the fine mesh as the verification reference (D-034). The reasons are:
\begin{itemize}
\item its error is small and conservative for thermal stress and growth (maximum solid temperature +4.4~K, volume mean +3.1~K above the extrapolated values);
\item the mesh error is about one quarter of the 19~K difference between the CFD and the analytical model, so the dominant thermal uncertainty is the
      heat-transfer and friction modelling, which a finer mesh does not reduce;
\item the fine mesh costs 3.4 times the run time and 4 times the storage;
\item the medium solid fits a 1:1 structural mesh under the 128,000-node Student limit, whereas the fine solid would not.
\end{itemize}
The structural consequence of the mesh error was carried forward as an uncertainty band: LC2 stress $-$10.1 / +2.8~MPa, local peak
$-$13.7 / +4.3~MPa and $\lambda_1$ +1.8 / $-$0.5\,\% (Chapter~\ref{ch:unc}).
